\section{What this guide certifies}
The completed fixed example is a bridge between exact arithmetic and ordinary
complex linear algebra. It certifies a twelve-dimensional Fourier matrix, a
Ramanujan projector, two normalized states with the same initial probabilities,
and their opposite outcomes after a basis change. This guide maps those claims to
actual Lean declarations rather than treating numerical agreement as a proof.

The five-module formal package contains \textbf{96 named mathematical theorems}. Its current verification rebuilt the source, inspected the axiom dependencies of every named theorem, and replayed all five modules. Historical proof records remain separate from this fresh execution.

\begin{center}
\begin{tabular}{lr}
\toprule
Production module & Named theorems\\
\midrule
\texttt{QuantumArithmeticBridge} & 22\\
\texttt{FiniteFourier} & 17\\
\texttt{SemanticBridge} & 30\\
\texttt{Conventions} & 17\\
\texttt{Publication} & 10\\
\midrule
Total & 96\\
\bottomrule
\end{tabular}
\end{center}

These counts include supporting lemmas. They count named source declarations,
not new discoveries, compiler jobs, or declarations generated internally by Lean.
The namespace is \texttt{Komorebi.QuantumArithmetic} throughout.

\section{The typed mathematical target}
The complex matrices act on vectors indexed by \(\operatorname{Fin}(12)\).
The Fourier convention is
\[
 F_{nk}=\frac{1}{\sqrt{12}}\exp\!\left(\frac{2\pi i nk}{12}\right),
 \qquad n,k\in\{0,\ldots,11\}.
\]
The adjoint \(F^*\) therefore carries the negative exponential sign. The
selector \(D\) is diagonal with value one precisely at the residues
\(U=\{1,5,7,11\}\). The transported operator is \(P=FDF^*\).

Write
\[
 c=(4,0,2,0,-2,0,-4,0,-2,0,2,0),\qquad
 \psi=\frac{c}{\sqrt{48}},\qquad \chi=\frac{|c|}{\sqrt{48}}.
\]
Both vectors have norm one. Their initial probability histograms are identical,
but their inner product is zero. The vector \(u\) has entry \(1/2\) on \(U\)
and zero elsewhere. Lean proves the actual complex identities \(Fu=\psi\) and
\(F^*\psi=u\), as well as transformed coprime probability one for \(\psi\)
and zero for \(\chi\).

The statements are finite and exact. No parameterized arbitrary-modulus theorem,
general circuit construction, physical phase-removal process, or quantum speedup
is included in these formal claims.

\section{Revision reconciliation}
Two retrieved archives carry version \texttt{0.2.0}, but they do not have the same
active source. The earlier four-module checkpoint has 59 mathematical theorems
and four separately checked refutation controls. The later five-module checkpoint
has 96 mathematical theorems and three rejected control files. Its proof receipt
is the later \texttt{verification/final/REPORT.json}.

The original 22 finite-table statements are retained inside both checkpoints.
They are not counted a second time. Likewise, \(59+96\) is not a meaningful
theorem total: these are overlapping revisions of one formal development.

The earlier proof boundary disclosed that explicit projector transport and
normalized-state matrix actions had not yet been supplied. The later source
supplies both, connects the Fourier matrix to the usual complex exponential, and
proves a histogram nonreconstruction theorem. The earlier boundary remains a
record of its own revision, rather than a gap list for the current publication.

The entire later checkpoint is preserved under \texttt{lean/checked\_checkpoint}.
The source files are not edited to match this guide.

\section{Architecture of the proof}
The proof crosses two representations. First, rational coordinates support exact
finite calculations. Second, a semantics map transfers those calculations to
ordinary complex matrices and vectors.

\begin{center}
\begin{tabular}{p{0.27\linewidth}p{0.64\linewidth}}
\toprule
Layer & Role\\
\midrule
Rational tables & Signal, projector entries, filtering, and energies.\\
Exact coordinates & Algebra in the basis \(1,\sqrt3,i,i\sqrt3\).\\
Finite matrices & All entries indexed by \(\operatorname{Fin}(12)\), with exact finite sums.\\
Complex semantics & Preservation of multiplication, adjoint, matrix action, and squared norm.\\
Conventions & Identification of roots and normalization with standard exponentials and square roots.\\
Publication theorem & Bundling of the actual complex identities and reconstruction obstructions.\\
\bottomrule
\end{tabular}
\end{center}

The active import chain is \texttt{QuantumArithmeticBridge} into
\texttt{FiniteFourier}, then \texttt{SemanticBridge}, then \texttt{Conventions},
then \texttt{Publication}. Supporting Mathlib imports enter at the layer that
uses them. There is no circular theorem dependency.

\section{The rational signal and its filter}
The first layer uses literal rational lists. The chosen perturbation is
\(r_n=3\cos(\pi n/2)+2\), whose displayed rational entries can be certified
without numerical trigonometric evaluation. Its sum with \(c\) is the test signal.

\Needspace{16\baselineskip}
\noindent\textbf{Literal source.} \texttt{\detokenize{QuantumArithmeticBridge.lean}}, physical lines 51--61.

\begin{Verbatim}[breaklines=true,breakanywhere=true,fontsize=\footnotesize,samepage=true]
/-- Ramanujan sum c_12(n), n = 0,...,11. -/
def ramanujan12 : QVector :=
  [4, 0, 2, 0, -2, 0, -4, 0, -2, 0, 2, 0]

/-- A deliberately off-band perturbation: 3 cos(pi n/2) + 2. -/
def noise12 : QVector :=
  [5, 2, -1, 2, 5, 2, -1, 2, 5, 2, -1, 2]

/-- Signal used in the essay's filter demonstration. -/
def testSignal12 : QVector :=
  [9, 2, 1, 2, 3, 2, -5, 2, 3, 2, 1, 2]
\end{Verbatim}

The cyclic matrix has entry \(c_{(n-m)\bmod12}/12\). The original layer proves
its symmetry, idempotence, preservation of \(c\), rejection of the chosen
perturbation, and exact recovery from the displayed contaminated signal.
The exact energy contract is \(150=48+102\).

\Needspace{6\baselineskip}
\noindent\textbf{Literal source.} \texttt{\detokenize{QuantumArithmeticBridge.lean}}, physical lines 92--126.

\begin{Verbatim}[breaklines=true,breakanywhere=true,fontsize=\footnotesize]
theorem projector12_idempotent :
    matMulN 12 projector12 projector12 = projector12 := by
  decide +kernel

/-- The arithmetic signal lies in the selected subspace. -/
theorem projector12_fixes_ramanujan :
    matVec projector12 ramanujan12 = ramanujan12 := by
  decide +kernel

/-- The chosen constant-plus-frequency-3 perturbation lies in the complementary subspace. -/
theorem projector12_kills_noise :
    matVec projector12 noise12 = List.replicate 12 0 := by
  decide +kernel

/-- Therefore the exact filter recovers the arithmetic component. -/
theorem projector12_recovers_test_signal :
    matVec projector12 testSignal12 = ramanujan12 := by
  decide +kernel

/-- The arithmetic amplitude vector has squared norm 48. -/
theorem ramanujan12_energy : sqEnergy ramanujan12 = 48 := by
  decide +kernel

/-- The perturbation has squared norm 102. -/
theorem noise12_energy : sqEnergy noise12 = 102 := by
  decide +kernel

/-- The full signal has squared norm 150. -/
theorem testSignal12_energy : sqEnergy testSignal12 = 150 := by
  decide +kernel

/-- Exact orthogonal energy decomposition 150 = 48 + 102. -/
theorem energy_decomposition :
    sqEnergy testSignal12 = sqEnergy ramanujan12 + sqEnergy noise12 := by
  decide +kernel
\end{Verbatim}

The list operations use finite lists, including zero defaults and zip-based dot
products. These facts are therefore claims about the supplied twelve-entry
objects. They should not be silently interpreted as universal matrix theorems
for arbitrary ragged lists. The later \(\operatorname{Fin}(12)\) layer supplies
the typed matrix identities used in the publication.

\section{The exact phase arithmetic}
The structure \texttt{Cyclo12} represents
\[
 (a,b,c,d)\longmapsto a+b\sqrt3+i(c+d\sqrt3),\qquad a,b,c,d\in\mathbb Q.
\]
Its multiplication explicitly incorporates \((\sqrt3)^2=3\) and \(i^2=-1\).
It is a concrete arithmetic representation. The development does not need or
install a global field instance to certify the finite example.

\Needspace{17\baselineskip}
\noindent\textbf{Literal source.} \texttt{\detokenize{QuantumArithmeticBridge.lean}}, physical lines 130--141.

\begin{Verbatim}[breaklines=true,breakanywhere=true,fontsize=\footnotesize,samepage=true]
/--
An exact coordinate model for Q(sqrt(3), i).  The tuple (a,b,c,d) denotes
  a + b*sqrt(3) + i*(c + d*sqrt(3)).
Only the concrete arithmetic needed for the q=12 certificate is used below;
we do not install a global field instance.
-/
structure Cyclo12 where
  a : ℚ
  b : ℚ
  c : ℚ
  d : ℚ
  deriving DecidableEq, Repr
\end{Verbatim}

\Needspace{6\baselineskip}
\noindent\textbf{Literal source.} \texttt{\detokenize{QuantumArithmeticBridge.lean}}, physical lines 155--173.

\begin{Verbatim}[breaklines=true,breakanywhere=true,fontsize=\footnotesize]
/-- Multiplication in Q(sqrt(3), i). -/
def cMul (x y : Cyclo12) : Cyclo12 :=
  ⟨x.a*y.a + 3*x.b*y.b - x.c*y.c - 3*x.d*y.d,
   x.a*y.b + x.b*y.a - x.c*y.d - x.d*y.c,
   x.a*y.c + 3*x.b*y.d + x.c*y.a + 3*x.d*y.b,
   x.a*y.d + x.b*y.c + x.c*y.b + x.d*y.a⟩

def cConj (x : Cyclo12) : Cyclo12 := ⟨x.a, x.b, -x.c, -x.d⟩

def cPow (x : Cyclo12) : Nat → Cyclo12
  | 0 => cOne
  | n + 1 => cMul (cPow x n) x

def cSum : List Cyclo12 → Cyclo12
  | [] => cZero
  | x :: xs => cAdd x (cSum xs)

/-- zeta = exp(2*pi*i/12) = sqrt(3)/2 + i/2. -/
def zeta12 : Cyclo12 := ⟨0, (1/2 : ℚ), (1/2 : ℚ), 0⟩
\end{Verbatim}

The phase \(\zeta=(0,1/2,1/2,0)\) has complex meaning
\(\sqrt3/2+i/2\). Exact finite certificates establish its twelfth power, inverse,
and absence of a smaller positive power below twelve. Later the convention
module identifies it with \(\exp(\pi i/6)\).

The Ramanujan vector is obtained by summing powers of this phase over
\(1,5,7,11\). The negative-sign unnormalized DFT has value 12 at these four
frequencies and zero elsewhere. The absolute-value vector has Fourier support
on the even frequencies instead. These are exact phase calculations over
rational coordinates.

\section{Finite indices and normalized Fourier certificates}
The finite layer uses \texttt{Index := Fin 12}. Its normalized Fourier scale is
represented exactly as \(\sqrt3/6=1/\sqrt{12}\), and its state scale as
\(\sqrt3/12=1/\sqrt{48}\). The use of a finite index type makes every row and
column well defined without list-length side conditions.

\Needspace{6\baselineskip}
\noindent\textbf{Literal source.} \texttt{\detokenize{FiniteFourier.lean}}, physical lines 58--93.

\begin{Verbatim}[breaklines=true,breakanywhere=true,fontsize=\footnotesize]
/-- The stored arithmetic selector is exactly the coprimality predicate. -/
theorem coprime_residues_exact :
    (List.range 12).filter (fun a => Nat.gcd a 12 == 1) = unitResidues12 := by
  decide +kernel

theorem root_exact_order :
    ∀ n : Fin 12, 0 < n.val → cPow zeta12 n.val ≠ cOne := by
  decide +kernel

theorem normalized_fourier_left :
    ∀ i j, kMatMul (kAdjoint exactFourier) exactFourier i j = kIdentity i j := by
  decide +kernel

theorem normalized_fourier_right :
    ∀ i j, kMatMul exactFourier (kAdjoint exactFourier) i j = kIdentity i j := by
  decide +kernel

def selectedKernel : KMatrix := fun i j =>
  kSum (fun k => if Nat.gcd k.val 12 = 1 then
    cMul (exactFourier i k) (cConj (exactFourier j k)) else cZero)

theorem exact_projector_transport :
    ∀ i j, selectedKernel i j = exactProjector i j := by
  decide +kernel

theorem exact_projector_matches_rational_table :
    ∀ i j, exactProjector i j = cRat ((projector12.getD i.val []).getD j.val 0) := by
  decide +kernel

theorem exact_state_preparation :
    ∀ i, kMatVec exactFourier exactUnitState i = exactPsi i := by
  decide +kernel

theorem exact_state_recovery :
    ∀ i, kMatVec (kAdjoint exactFourier) exactPsi i = exactUnitState i := by
  decide +kernel
\end{Verbatim}

The two Fourier identities above quantify over all \(12\times12\) entries.
Kernel reduction proves these concrete universal statements over a finite type.
The selected-kernel equality and state actions are also complete finite
certificates, rather than floating-point matrix comparisons.

\section{Why the coordinate certificates mean complex identities}
The semantics map is explicitly defined with the positive real square root.
Soundness for addition, multiplication, conjugation, powers, and finite sums
prevents the exact coordinate layer from being a detached calculation with an
unproved interpretation.

\Needspace{6\baselineskip}
\noindent\textbf{Literal source.} \texttt{\detokenize{SemanticBridge.lean}}, physical lines 14--56.

\begin{Verbatim}[breaklines=true,breakanywhere=true,fontsize=\footnotesize]
/-- The exact coordinate arithmetic is interpreted in the ordinary complex numbers. -/
noncomputable def interpret (x : Cyclo12) : ℂ :=
  ⟨(x.a : ℝ) + (x.b : ℝ) * Real.sqrt 3,
   (x.c : ℝ) + (x.d : ℝ) * Real.sqrt 3⟩

@[simp] theorem interpret_zero : interpret cZero = 0 := by
  apply Complex.ext <;> simp [interpret, cZero]

@[simp] theorem interpret_one : interpret cOne = 1 := by
  apply Complex.ext <;> simp [interpret, cOne]

@[simp] theorem interpret_rat (r : ℚ) : interpret (cRat r) = (r : ℂ) := by
  apply Complex.ext <;> simp [interpret, cRat]

@[simp] theorem interpret_add (x y : Cyclo12) :
    interpret (cAdd x y) = interpret x + interpret y := by
  apply Complex.ext <;> simp [interpret, cAdd] <;> ring

@[simp] theorem interpret_mul (x y : Cyclo12) :
    interpret (cMul x y) = interpret x * interpret y := by
  apply Complex.ext <;> simp [interpret, cMul, Complex.mul_re, Complex.mul_im] <;>
    ring_nf <;> norm_num [Real.sq_sqrt]

@[simp] theorem interpret_conj (x : Cyclo12) :
    interpret (cConj x) = star (interpret x) := by
  apply Complex.ext <;> simp [interpret, cConj]
  ring

@[simp] theorem interpret_pow (x : Cyclo12) (n : Nat) :
    interpret (cPow x n) = interpret x ^ n := by
  induction n with
  | zero => simp [cPow]
  | succ n ih => simp [cPow, ih, pow_succ]

@[simp] theorem interpret_sum (xs : List Cyclo12) :
    interpret (cSum xs) = (xs.map interpret).sum := by
  induction xs with
  | nil => simp [cSum]
  | cons x xs ih => simp [cSum, ih]

@[simp] theorem interpret_kSum (f : Index → Cyclo12) :
    interpret (kSum f) = ∑ i, interpret (f i) := by
  simp only [kSum, interpret_sum, List.map_ofFn, List.sum_ofFn]
\end{Verbatim}

The multiplication proof uses complex extensionality, real polynomial
normalization, and the square-root identity. It proves preservation for arbitrary
coordinates, not merely for the handful of roots that occur in the example.
The power and finite-sum proofs use induction.

Injectivity of this interpretation is not needed to transfer an equality: if
two coordinate expressions are proved equal, applying the interpretation gives
equal complex expressions. No claim of a globally packaged field equivalence is
required by the theorem chain.

\section{Matrix transport and actual unitarity}
The semantics layer proves that coordinate matrix multiplication, adjoint,
identity, and matrix-vector action have their ordinary complex meanings.
The matrices \(F,P,D\) and the normalized states are then defined directly
in that complex representation.

\Needspace{6\baselineskip}
\noindent\textbf{Literal source.} \texttt{\detokenize{SemanticBridge.lean}}, physical lines 65--91.

\begin{Verbatim}[breaklines=true,breakanywhere=true,fontsize=\footnotesize]
@[simp] theorem complexMatrix_mul (A B : KMatrix) :
    complexMatrix (kMatMul A B) = complexMatrix A * complexMatrix B := by
  ext i j
  simp [complexMatrix, kMatMul, Matrix.mul_apply]

@[simp] theorem complexMatrix_adjoint (A : KMatrix) :
    complexMatrix (kAdjoint A) = (complexMatrix A)ᴴ := by
  ext i j
  simp [complexMatrix, kAdjoint, Matrix.conjTranspose_apply]

@[simp] theorem complexMatrix_identity :
    complexMatrix kIdentity = 1 := by
  ext i j
  by_cases h : i = j <;> simp [complexMatrix, kIdentity, h, Matrix.one_apply]

@[simp] theorem complexVector_mul (A : KMatrix) (v : KVector) :
    complexVector (kMatVec A v) = (complexMatrix A).mulVec (complexVector v) := by
  ext i
  simp [complexVector, complexMatrix, kMatVec, Matrix.mulVec, dotProduct]

noncomputable def F : Matrix Index Index ℂ := complexMatrix exactFourier
noncomputable def P : Matrix Index Index ℂ := complexMatrix exactProjector
noncomputable def D : Matrix Index Index ℂ :=
  Matrix.diagonal (fun i => if Nat.gcd i.val 12 = 1 then 1 else 0)
noncomputable def psi : Index → ℂ := complexVector exactPsi
noncomputable def chi : Index → ℂ := complexVector exactChi
noncomputable def unitState : Index → ℂ := complexVector exactUnitState
\end{Verbatim}

Consequently, the finite certificates produce actual complex unitarity and
projector transport. The following excerpt contains the full proofs of both
unitary identities, the cross-representation projector identity, and the two
state-action identities.

\Needspace{6\baselineskip}
\noindent\textbf{Literal source.} \texttt{\detokenize{SemanticBridge.lean}}, physical lines 93--120.

\begin{Verbatim}[breaklines=true,breakanywhere=true,fontsize=\footnotesize]
/-- Actual complex unitarity, not an assertion only about the coordinate checker. -/
theorem fourier_unitary_left : Fᴴ * F = 1 := by
  have h : kMatMul (kAdjoint exactFourier) exactFourier = kIdentity :=
    funext (fun i => funext (normalized_fourier_left i))
  simpa [F] using congrArg complexMatrix h

theorem fourier_unitary_right : F * Fᴴ = 1 := by
  have h : kMatMul exactFourier (kAdjoint exactFourier) = kIdentity :=
    funext (fun i => funext (normalized_fourier_right i))
  simpa [F] using congrArg complexMatrix h

theorem projector_transport : F * D * Fᴴ = P := by
  ext i j
  have h := congrArg interpret (exact_projector_transport i j)
  simp only [selectedKernel, interpret_kSum] at h
  rw [Matrix.mul_apply]
  simpa [D, P, F, complexMatrix, Matrix.mul_diagonal, Matrix.conjTranspose_apply,
    apply_ite interpret, mul_ite] using h

theorem state_preparation : F.mulVec unitState = psi := by
  have h : kMatVec exactFourier exactUnitState = exactPsi :=
    funext exact_state_preparation
  simpa [F, unitState, psi] using congrArg complexVector h

theorem state_recovery : Fᴴ.mulVec psi = unitState := by
  have h : kMatVec (kAdjoint exactFourier) exactPsi = exactUnitState :=
    funext exact_state_recovery
  simpa [F, unitState, psi] using congrArg complexVector h
\end{Verbatim}

This is the key strengthening over the earlier checkpoint. The rational matrix,
the exact coordinate kernel, and the complex matrix \(FDF^*\) are explicitly
connected. State preparation and recovery are equalities of actual complex
vectors. The theorem name \texttt{state\_preparation} describes a matrix identity,
not an efficient hardware preparation algorithm.

\section{The exponential sign and square-root normalization}
The usual Fourier matrix is not left as an informal interpretation. The
convention module proves that the coordinate root equals the positive
exponential root, including the selected square-root branch.

\Needspace{20\baselineskip}
\noindent\textbf{Literal source.} \texttt{\detokenize{Conventions.lean}}, physical lines 10--24.

\begin{Verbatim}[breaklines=true,breakanywhere=true,fontsize=\footnotesize,samepage=true]
/-- The coordinate root is exactly the usual exponential twelfth root. -/
theorem root_is_exponential :
    interpret zeta12 = Complex.exp (((Real.pi / 6 : ℝ) : ℂ) * Complex.I) := by
  rw [Complex.exp_mul_I]
  rw [← Complex.ofReal_cos, ← Complex.ofReal_sin, Real.cos_pi_div_six, Real.sin_pi_div_six]
  apply Complex.ext <;> simp [interpret, zeta12]
  ring

theorem root_power_twelve : interpret zeta12 ^ 12 = 1 := by
  simpa using congrArg interpret zeta12_order_certificate

theorem rootPower_interpret (n : Nat) :
    interpret (rootPower n) = interpret zeta12 ^ n := by
  simp only [rootPower, interpret_pow]
  exact (pow_eq_pow_mod n root_power_twelve).symm
\end{Verbatim}

Its normalization lemmas prove \(\sqrt{12}=2\sqrt3\) and
\(\sqrt{48}=4\sqrt3\), then identify the exact scales with their ordinary
reciprocals. The matrix entry formula and state formula below are therefore
literal equalities in \(\mathbb C\).

\Needspace{18\baselineskip}
\noindent\textbf{Literal source.} \texttt{\detokenize{Conventions.lean}}, physical lines 42--54.

\begin{Verbatim}[breaklines=true,breakanywhere=true,fontsize=\footnotesize,samepage=true]
theorem fourierScale_interpret :
    interpret fourierScale = ((1 / Real.sqrt 12 : ℝ) : ℂ) := by
  apply Complex.ext <;> simp [interpret, fourierScale, sqrt_twelve]
  have h : Real.sqrt 3 ≠ 0 := by positivity
  field_simp
  nlinarith [Real.sq_sqrt (show (0 : ℝ) ≤ 3 by norm_num)]

theorem stateScale_interpret :
    interpret stateScale = ((1 / Real.sqrt 48 : ℝ) : ℂ) := by
  apply Complex.ext <;> simp [interpret, stateScale, sqrt_fortyeight]
  have h : Real.sqrt 3 ≠ 0 := by positivity
  field_simp
  nlinarith [Real.sq_sqrt (show (0 : ℝ) ≤ 3 by norm_num)]
\end{Verbatim}

\Needspace{15\baselineskip}
\noindent\textbf{Literal source.} \texttt{\detokenize{Conventions.lean}}, physical lines 56--65.

\begin{Verbatim}[breaklines=true,breakanywhere=true,fontsize=\footnotesize,samepage=true]
/-- The matrix checked in the certificate is the standard positive sign Fourier matrix. -/
theorem fourier_entry_formula (i j : Index) :
    F i j = ((1 / Real.sqrt 12 : ℝ) : ℂ) *
      Complex.exp (2 * (Real.pi : ℂ) * Complex.I * (i.val : ℂ) * (j.val : ℂ) / 12) := by
  simp only [F, complexMatrix, exactFourier, interpret_mul, fourierScale_interpret,
    rootPower_interpret, root_is_exponential]
  rw [← Complex.exp_nat_mul]
  congr 2
  push_cast
  ring
\end{Verbatim}

\Needspace{21\baselineskip}
\noindent\textbf{Literal source.} \texttt{\detokenize{Conventions.lean}}, physical lines 67--82.

\begin{Verbatim}[breaklines=true,breakanywhere=true,fontsize=\footnotesize,samepage=true]
theorem psi_entry_formula (i : Index) :
    psi i = ((1 / Real.sqrt 48 : ℝ) : ℂ) * (ramanujan12.getD i.val 0 : ℂ) := by
  simp [psi, complexVector, exactPsi, stateScale_interpret]

theorem chi_entry_formula (i : Index) :
    chi i = ((1 / Real.sqrt 48 : ℝ) : ℂ) * (absRamanujan12.getD i.val 0 : ℂ) := by
  simp [chi, complexVector, exactChi, stateScale_interpret]

theorem unitState_entry_formula (i : Index) :
    unitState i = if Nat.gcd i.val 12 = 1 then (1/2 : ℂ) else 0 := by
  by_cases h : Nat.gcd i.val 12 = 1 <;>
    simp [unitState, complexVector, exactUnitState, h]

theorem projector_entry_formula (i j : Index) :
    P i j = (ramanujan12.getD ((i.val + 12 - j.val) % 12) 0 : ℂ) / 12 := by
  simp [P, complexMatrix, exactProjector]
\end{Verbatim}

The convention is positive sign for \(F\) and negative sign for \(F^*\). The
earlier unnormalized \texttt{dft12} uses the negative sign. The normalization
and adjoint are the precise link between those notations.

\section{Hermitian projection and phase-sensitive states}
The selector is Hermitian and idempotent. Transport through a unitary matrix
then proves the corresponding facts for \(P\). They are ordinary matrix
identities, not claims that every rational list operation is an orthogonal
projection.

\Needspace{6\baselineskip}
\noindent\textbf{Literal source.} \texttt{\detokenize{Conventions.lean}}, physical lines 84--108.

\begin{Verbatim}[breaklines=true,breakanywhere=true,fontsize=\footnotesize]
theorem selector_idempotent : D * D = D := by
  unfold D
  rw [Matrix.diagonal_mul_diagonal]
  congr 1
  funext i
  by_cases h : Nat.gcd i.val 12 = 1 <;> simp [h]

theorem selector_selfAdjoint : Dᴴ = D := by
  ext i j
  by_cases h : i = j
  · subst j
    simp [D, Matrix.conjTranspose_apply, apply_ite star]
  · simp [D, Matrix.conjTranspose_apply, h, Ne.symm h]

theorem projector_idempotent_complex : P * P = P := by
  rw [← projector_transport]
  calc
    (F * D * Fᴴ) * (F * D * Fᴴ) = F * D * (Fᴴ * F) * D * Fᴴ := by
      simp only [Matrix.mul_assoc]
    _ = F * D * Fᴴ := by rw [fourier_unitary_left, Matrix.mul_one, Matrix.mul_assoc F D D,
      selector_idempotent]

theorem projector_selfAdjoint_complex : Pᴴ = P := by
  rw [← projector_transport]
  simp [Matrix.conjTranspose_mul, selector_selfAdjoint, Matrix.mul_assoc]
\end{Verbatim}

The vectors \(\psi\) and \(\chi\) are separately defined states. One is the
signed arithmetic amplitude vector and the other its entrywise absolute value.
They have the same squared amplitudes at each initial index but retain different
relative phase. The development neither claims nor requires a universal quantum
operation implementing the nonlinear map \(v\mapsto |v|\).

\Needspace{6\baselineskip}
\noindent\textbf{Literal source.} \texttt{\detokenize{SemanticBridge.lean}}, physical lines 145--174.

\begin{Verbatim}[breaklines=true,breakanywhere=true,fontsize=\footnotesize]
theorem psi_normalized : totalProbability psi = 1 := by
  have h := congrArg interpret exact_psi_normalized
  simpa [psi] using congrArg Complex.re h

theorem chi_normalized : totalProbability chi = 1 := by
  have h := congrArg interpret exact_chi_normalized
  simpa [chi] using congrArg Complex.re h

theorem unitState_normalized : totalProbability unitState = 1 := by
  have h := congrArg interpret exact_unit_normalized
  simpa [unitState] using congrArg Complex.re h

theorem same_initial_probabilities (i : Index) :
    Complex.normSq (psi i) = Complex.normSq (chi i) := by
  have h := congrArg interpret (exact_same_initial_probabilities i)
  simpa [psi, chi, complexVector] using congrArg Complex.re h

theorem states_orthogonal : ∑ i, star (psi i) * chi i = 0 := by
  have h := congrArg interpret exact_orthogonal_states
  simpa [psi, chi, complexVector] using h

theorem signed_state_recovers_arithmetic :
    coprimeProbability (Fᴴ.mulVec psi) = 1 := by
  have h := congrArg interpret exact_coprime_mass_psi
  simpa [F, psi] using congrArg Complex.re h

theorem unsigned_state_erases_arithmetic :
    coprimeProbability (Fᴴ.mulVec chi) = 0 := by
  have h := congrArg interpret exact_coprime_mass_chi
  simpa [F, chi] using congrArg Complex.re h
\end{Verbatim}

The transformed probability statement is exact. The measurement event is the
coprime-frequency subset. After applying \(F^*\), its probability is one for
\(\psi\) and zero for \(\chi\). This basis change reveals information that their
shared initial histogram does not contain.

\section{Two exact information-loss statements}
The formal development isolates distinct obstructions. First, the initial
histogram cannot be inverted on both states by any reconstruction function.
Second, the projector itself is noninjective because it annihilates the nonzero
normalized state \(\chi\).

\Needspace{17\baselineskip}
\noindent\textbf{Literal source.} \texttt{\detokenize{Publication.lean}}, physical lines 39--50.

\begin{Verbatim}[breaklines=true,breakanywhere=true,fontsize=\footnotesize,samepage=true]
/-- No state reconstruction function can invert the initial probability histogram on both states. -/
theorem no_reconstruction_from_initial_probabilities :
    ¬ ∃ R : (Index → ℝ) → (Index → ℂ),
      R (fun i => Complex.normSq (psi i)) = psi ∧
      R (fun i => Complex.normSq (chi i)) = chi := by
  rintro ⟨R, hpsi, hchi⟩
  have heq : (fun i => Complex.normSq (psi i)) =
      (fun i => Complex.normSq (chi i)) := funext same_initial_probabilities
  have hstate : psi = chi := by rw [← hpsi, ← hchi, heq]
  have h := signed_state_recovers_arithmetic
  rw [hstate, unsigned_state_erases_arithmetic] at h
  norm_num at h
\end{Verbatim}

\Needspace{11\baselineskip}
\noindent\textbf{Literal source.} \texttt{\detokenize{Publication.lean}}, physical lines 52--57.

\begin{Verbatim}[breaklines=true,breakanywhere=true,fontsize=\footnotesize,samepage=true]
/-- Exact probabilities in the recovered basis, connected to the normalized state. -/
theorem recovered_probability_at (i : Index) :
    Complex.normSq (Fᴴ.mulVec psi i) =
      if Nat.gcd i.val 12 = 1 then (1/4 : ℝ) else 0 := by
  rw [state_recovery, unitState_entry_formula]
  by_cases h : Nat.gcd i.val 12 = 1 <;> norm_num [h, Complex.normSq_apply]
\end{Verbatim}

\Needspace{12\baselineskip}
\noindent\textbf{Literal source.} \texttt{\detokenize{Publication.lean}}, physical lines 59--65.

\begin{Verbatim}[breaklines=true,breakanywhere=true,fontsize=\footnotesize,samepage=true]
/-- The operator is not injective. Its nonzero normalized kernel state is explicit. -/
theorem projector_not_injective : ¬ Function.Injective P.mulVec := by
  intro h
  have hz : chi = 0 := h (by simp [projector_kills_chi])
  have hp := chi_normalized
  rw [hz] at hp
  norm_num [totalProbability] at hp
\end{Verbatim}

The nonreconstruction theorem quantifies over every function from real-valued
histograms to complex-valued vectors. It uses equality of the two histograms and
the incompatible transformed probabilities. This is a mathematical obstruction
to a specified inverse, not a statistical claim about estimating a state from
many measurements in several bases.

The traditional paper also proves that the two states cannot both be recovered
up to an overall phase from this histogram. That strengthening is a prose
corollary of their orthogonality and normalization. It is not separately exported
as a quotient-by-global-phase Lean theorem.

Projection becomes reconstructive when its residual is retained. The next
identity holds for every twelve-dimensional complex vector, including vectors
outside the arithmetic signal subspace.

\Needspace{10\baselineskip}
\noindent\textbf{Literal source.} \texttt{\detokenize{Publication.lean}}, physical lines 33--37.

\begin{Verbatim}[breaklines=true,breakanywhere=true,fontsize=\footnotesize,samepage=true]
/-- Every signal is reconstructible from the retained and discarded components together. -/
theorem projector_residual_reconstructs (v : Index → ℂ) :
    P.mulVec v + (1 - P).mulVec v = v := by
  rw [Matrix.sub_mulVec, Matrix.one_mulVec]
  exact add_sub_cancel _ _
\end{Verbatim}

Its proof is algebraic. The identity \(Pv+(I-P)v=v\) itself does not require
idempotence. Hermitian idempotence supplies the stronger interpretation as
orthogonal retained and discarded components. This separates the reconstruction
identity from the additional geometric structure.

\section{The publication theorem in full}
The final declaration bundles the formal core of the worked example. Every
conjunct is supported by a previously proved declaration rather than by a new
numerical test or an implicit physical assumption.

\Needspace{19\baselineskip}
\noindent\textbf{Literal source.} \texttt{\detokenize{Publication.lean}}, physical lines 67--80.

\begin{Verbatim}[breaklines=true,breakanywhere=true,fontsize=\footnotesize,samepage=true]
/-- Publication theorem gathering the complex identities stated in the worked example. -/
theorem publication_bridge :
    Fᴴ * F = 1 ∧ F * Fᴴ = 1 ∧ F * D * Fᴴ = P ∧
    P * P = P ∧ Pᴴ = P ∧ F.mulVec unitState = psi ∧
    Fᴴ.mulVec psi = unitState ∧ totalProbability psi = 1 ∧
    totalProbability chi = 1 ∧
    (∀ i, Complex.normSq (psi i) = Complex.normSq (chi i)) ∧
    (∑ i, star (psi i) * chi i) = 0 ∧
    coprimeProbability (Fᴴ.mulVec psi) = 1 ∧
    coprimeProbability (Fᴴ.mulVec chi) = 0 := by
  exact ⟨fourier_unitary_left, fourier_unitary_right, projector_transport,
    projector_idempotent_complex, projector_selfAdjoint_complex, state_preparation,
    state_recovery, psi_normalized, chi_normalized, same_initial_probabilities,
    states_orthogonal, signed_state_recovers_arithmetic, unsigned_state_erases_arithmetic⟩
\end{Verbatim}

The bundled theorem has a fixed dimension and explicitly named operators and
states. It establishes the bridge used in the paper. Broader arbitrary-modulus
results remain separate prose theorems, with their own proofs and evidence.

The manuscript's QA18 additionally includes filtering and energy, normalization
of \(u\), the projector actions on both states, and noninjectivity. Those clauses
use separately exported declarations, including
\texttt{projector12\_recovers\_test\_signal}, \texttt{energy\_decomposition},
\texttt{unitState\_normalized}, \texttt{projector\_fixes\_psi},
\texttt{projector\_kills\_chi}, and \texttt{projector\_not\_injective}.
The complete machine concordance lists these supporting results. The manuscript
bundle is not identified with a literal statement that omits its extra clauses.

\section{Manuscript-to-Lean concordance}
The main paper labels its results QA01 through QA20. The table below distinguishes
the exact fixed-example formal coverage from the more general mathematical
development. ``Not formalized here'' is a scope label, not a claim that the prose
result is false.

\small
\begin{longtable}{p{0.08\linewidth}p{0.26\linewidth}p{0.56\linewidth}}
\toprule
Claim & Subject & Formal coverage\\
\midrule
\endhead
QA01 & Primitive root & Exact order 12 and standard exponential root.\\
QA02 & Character orthogonality & Fixed-example content through both exact unitary identities; general theorem not separately exported.\\
QA03 & Fourier theory & Both ordinary complex unitary identities and standard entry formula; general Parseval theorem not separately exported.\\
QA04 & M\"obius formula & General modulus formula not formalized; exact integer table for 12 certified.\\
QA05 & Ramanujan support & Exact signed and unsigned twelve-frequency DFT tables and arithmetic energy 48.\\
QA06 & Divisor projectors & Primitive-order projector at 12 certified; the full divisor family not a general Lean family.\\
QA07 & Projector geometry & Hermitian idempotence, displayed energies, explicit noninjectivity, and all-vector residual reconstruction.\\
QA08 & Residue recombination & No theorem in this checkpoint.\\
QA09 & Polyphase aliasing & No theorem in this checkpoint.\\
QA10 & Arithmetic example & Phase sum connected to the exact rational sample table.\\
QA11 & Exact filtering & Displayed contaminated signal and orthogonal energy decomposition.\\
QA12 & State preparation & Ordinary complex normalized vector actions and recovery.\\
QA13 & Shared histogram & Entrywise probability equality and state orthogonality.\\
QA14 & Phase separation & Actual transformed coprime probabilities 1 versus 0.\\
QA15 & Reconstruction barrier & No vector-valued initial-histogram inverse works on both specified states; the global-phase corollary is prose only.\\
QA16 & Measurement law & Relevant projector algebra and state probabilities; no general conditional-state-update API.\\
QA17 & Arithmetic semantics & Soundness for scalar, matrix, vector, adjoint, and mass operations.\\
QA18 & Central bridge & Publication bundle plus separately exported filter, energy, normalization, projector-action, and noninjectivity results.\\
QA19 & Relative-phase family & Prose extension \(\sin^2(\theta/2)\); not an inherited Lean theorem.\\
QA20 & Dephased probability & Prose density-matrix calculation \(1/3\); not an inherited Lean theorem.\\
\bottomrule
\end{longtable}
\normalsize

The machine concordance includes all theorem statements verbatim, per-theorem
axiom records, source hashes, and exact physical line locators. Its mapping
records formal scope; it does not upgrade general manuscript claims merely
because a nearby finite special case is checked.

\section{Trust boundary and current verification}
The current replay uses Lean 4.34.0-rc2, compiler commit
\texttt{\detokenize{6a10ac8c22beadecabdbb0919c2b50214762f91d}}, and Mathlib revision
\texttt{\detokenize{aeb72f1e9981b6d35b426b16dca7759be3e7ba06}}.
It built the five source modules with warnings treated as errors, printed axiom
dependencies for all 96 named theorems, and replayed each module with
\texttt{leanchecker}. The compiler and imported dependency cache were not rebuilt
from their source. Kernel replay uses Lean's own kernel, not a second independent
proof checker.

The axiom union is exactly \texttt{propext}, \texttt{Classical.choice}, and
\texttt{Quot.sound}. This is the standard foundation used by the audited proof
terms, not an axiom-free claim. The active source contains no proof placeholders,
custom axiom declarations, native decision proofs, or unsafe declarations.
Historical failed or native-decision drafts are excluded from the active build.

Current source reconciliation checked all five source hashes and all eleven
archived check-log hashes against the archived report. All sixteen match.
The current inventory contains exactly 96 unique names and 96 archived axiom
records. All line numbers agree with the archived source index.

The pinned portable runtime was restored and verified against its combined archive hash. A fresh build, all 96 theorem axiom inspections, and all five kernel replays passed. The complete receipt and logs accompany the source package. Earlier transfer failures remain historical records.

\section{Controls, failed attempts, and reproducibility}
The latest archive rejected a deliberately false signal sign and a proof using
\texttt{sorry} under warnings as errors. Its false Fourier-mode control also
failed, but the compiler stopped at opaque \texttt{Decidable} reduction. That
failure alone is not a direct mathematical certificate that the proposition is
false. The positive \texttt{ramanujan\_dft\_table} theorem separately establishes
the correct value 12.

The control runner's broad diagnostic expression can match the word ``false''
in a filename. The guide does not interpret that match as a stronger arithmetic
test than the actual log provides. The earlier checkpoint's four explicit
refutation proofs are retained as earlier evidence and are not silently counted
as current production theorems.

A supplementary current check closes the diagnostic gap for the false-mode control. It derives the exact mode-one value and its negation from the proved Fourier table, replays the supplementary module through the kernel, and rejects a false rational-coordinate equality with an explicit false-value diagnostic. These two verification lemmas do not increase the monograph count of 96 theorems.

The supplied verifier creates a fresh source project, checks dependency revisions,
builds, prints all theorem axiom dependencies, replays all five modules, and runs
the controls. It refuses to overwrite an existing fresh build directory and
treats timeouts as failures. It makes no network requests or installations.
With the pinned runtime restored, the documented command is:

\begin{Verbatim}[breaklines=true,fontsize=\footnotesize]
python scripts/verify_lean.py \
  --toolchain /path/to/Komorebi_Portable_Lean_4.34.0-rc2/toolchain \
  --packages /path/to/Komorebi_Portable_Lean_4.34.0-rc2/workspace/.lake/packages \
  --out verification/reproduction_new
\end{Verbatim}
The runner was exercised under Linux x86-64. It uses POSIX process groups and
filesystem symlinks. Native Windows behavior is not tested. WSL2 with the pinned
Linux runtime is the documented route for a Windows machine.

Full proof source, manifests, archived build and replay logs, the static
concordance builder, and the current successful replay receipt are
included. The runtime itself remains an external dependency.

\section{Complete active theorem inventory}
All names in the following table have prefix
\texttt{Komorebi.QuantumArithmetic}. File and physical source line identify the
exact declaration. The full literal statements and axiom lists are also available
in \texttt{Lean\_Concordance.json} and \texttt{Lean\_Concordance.csv}.

\Needspace{10\baselineskip}
\subsection{\texttt{\detokenize{QuantumArithmeticBridge}}: 22 theorems}

\small\begin{longtable}{p{0.79\linewidth}r}
\toprule Declaration & Line\\\midrule\endhead
\texttt{\detokenize{ramanujan12_length}} & 80\\
\texttt{\detokenize{projector12_shape}} & 82\\
\texttt{\detokenize{projector12_symmetric}} & 88\\
\texttt{\detokenize{projector12_idempotent}} & 92\\
\texttt{\detokenize{projector12_fixes_ramanujan}} & 97\\
\texttt{\detokenize{projector12_kills_noise}} & 102\\
\texttt{\detokenize{projector12_recovers_test_signal}} & 107\\
\texttt{\detokenize{ramanujan12_energy}} & 112\\
\texttt{\detokenize{noise12_energy}} & 116\\
\texttt{\detokenize{testSignal12_energy}} & 120\\
\texttt{\detokenize{energy_decomposition}} & 124\\
\texttt{\detokenize{zeta12_order_certificate}} & 177\\
\texttt{\detokenize{zeta12_inverse_certificate}} & 180\\
\texttt{\detokenize{ramanujanC12_table}} & 190\\
\texttt{\detokenize{ramanujan_dft_table}} & 205\\
\texttt{\detokenize{abs_ramanujan_dft_table}} & 217\\
\texttt{\detokenize{same_computational_basis_probabilities}} & 231\\
\texttt{\detokenize{abs_ramanujan12_energy}} & 235\\
\texttt{\detokenize{recovered_mode_probability}} & 243\\
\texttt{\detokenize{recovered_total_probability}} & 248\\
\texttt{\detokenize{phase_stripping_coprime_probability}} & 253\\
\texttt{\detokenize{arithmetic_signal_quantum_bridge}} & 258\\
\bottomrule\end{longtable}\normalsize

\Needspace{10\baselineskip}
\subsection{\texttt{\detokenize{FiniteFourier}}: 17 theorems}

\small\begin{longtable}{p{0.79\linewidth}r}
\toprule Declaration & Line\\\midrule\endhead
\texttt{\detokenize{coprime_residues_exact}} & 59\\
\texttt{\detokenize{root_exact_order}} & 63\\
\texttt{\detokenize{normalized_fourier_left}} & 67\\
\texttt{\detokenize{normalized_fourier_right}} & 71\\
\texttt{\detokenize{exact_projector_transport}} & 79\\
\texttt{\detokenize{exact_projector_matches_rational_table}} & 83\\
\texttt{\detokenize{exact_state_preparation}} & 87\\
\texttt{\detokenize{exact_state_recovery}} & 91\\
\texttt{\detokenize{exact_unit_normalized}} & 95\\
\texttt{\detokenize{exact_psi_normalized}} & 98\\
\texttt{\detokenize{exact_chi_normalized}} & 101\\
\texttt{\detokenize{exact_same_initial_probabilities}} & 104\\
\texttt{\detokenize{exact_orthogonal_states}} & 108\\
\texttt{\detokenize{exact_coprime_mass_psi}} & 112\\
\texttt{\detokenize{exact_coprime_mass_chi}} & 116\\
\texttt{\detokenize{exact_projector_fixes_psi}} & 120\\
\texttt{\detokenize{exact_projector_kills_chi}} & 124\\
\bottomrule\end{longtable}\normalsize

\Needspace{10\baselineskip}
\subsection{\texttt{\detokenize{SemanticBridge}}: 30 theorems}

\small\begin{longtable}{p{0.79\linewidth}r}
\toprule Declaration & Line\\\midrule\endhead
\texttt{\detokenize{interpret_zero}} & 19\\
\texttt{\detokenize{interpret_one}} & 22\\
\texttt{\detokenize{interpret_rat}} & 25\\
\texttt{\detokenize{interpret_add}} & 28\\
\texttt{\detokenize{interpret_mul}} & 32\\
\texttt{\detokenize{interpret_conj}} & 37\\
\texttt{\detokenize{interpret_pow}} & 42\\
\texttt{\detokenize{interpret_sum}} & 48\\
\texttt{\detokenize{interpret_kSum}} & 54\\
\texttt{\detokenize{complexMatrix_mul}} & 65\\
\texttt{\detokenize{complexMatrix_adjoint}} & 70\\
\texttt{\detokenize{complexMatrix_identity}} & 75\\
\texttt{\detokenize{complexVector_mul}} & 80\\
\texttt{\detokenize{fourier_unitary_left}} & 94\\
\texttt{\detokenize{fourier_unitary_right}} & 99\\
\texttt{\detokenize{projector_transport}} & 104\\
\texttt{\detokenize{state_preparation}} & 112\\
\texttt{\detokenize{state_recovery}} & 117\\
\texttt{\detokenize{interpret_normSq}} & 122\\
\texttt{\detokenize{interpret_mass}} & 133\\
\texttt{\detokenize{interpret_coprimeMass}} & 137\\
\texttt{\detokenize{psi_normalized}} & 145\\
\texttt{\detokenize{chi_normalized}} & 149\\
\texttt{\detokenize{unitState_normalized}} & 153\\
\texttt{\detokenize{same_initial_probabilities}} & 157\\
\texttt{\detokenize{states_orthogonal}} & 162\\
\texttt{\detokenize{signed_state_recovers_arithmetic}} & 166\\
\texttt{\detokenize{unsigned_state_erases_arithmetic}} & 171\\
\texttt{\detokenize{projector_fixes_psi}} & 176\\
\texttt{\detokenize{projector_kills_chi}} & 181\\
\bottomrule\end{longtable}\normalsize

\Needspace{10\baselineskip}
\subsection{\texttt{\detokenize{Conventions}}: 17 theorems}

\small\begin{longtable}{p{0.79\linewidth}r}
\toprule Declaration & Line\\\midrule\endhead
\texttt{\detokenize{root_is_exponential}} & 11\\
\texttt{\detokenize{root_power_twelve}} & 18\\
\texttt{\detokenize{rootPower_interpret}} & 21\\
\texttt{\detokenize{sqrt_twelve}} & 26\\
\texttt{\detokenize{sqrt_fortyeight}} & 34\\
\texttt{\detokenize{fourierScale_interpret}} & 42\\
\texttt{\detokenize{stateScale_interpret}} & 49\\
\texttt{\detokenize{fourier_entry_formula}} & 57\\
\texttt{\detokenize{psi_entry_formula}} & 67\\
\texttt{\detokenize{chi_entry_formula}} & 71\\
\texttt{\detokenize{unitState_entry_formula}} & 75\\
\texttt{\detokenize{projector_entry_formula}} & 80\\
\texttt{\detokenize{selector_idempotent}} & 84\\
\texttt{\detokenize{selector_selfAdjoint}} & 91\\
\texttt{\detokenize{projector_idempotent_complex}} & 98\\
\texttt{\detokenize{projector_selfAdjoint_complex}} & 106\\
\texttt{\detokenize{complete_phase_separation}} & 111\\
\bottomrule\end{longtable}\normalsize

\Needspace{10\baselineskip}
\subsection{\texttt{\detokenize{Publication}}: 10 theorems}

\small\begin{longtable}{p{0.79\linewidth}r}
\toprule Declaration & Line\\\midrule\endhead
\texttt{\detokenize{exact_phase_signal_lookup}} & 9\\
\texttt{\detokenize{complex_ramanujan_signal}} & 14\\
\texttt{\detokenize{absolute_signal_lookup}} & 22\\
\texttt{\detokenize{retained_discarded_orthogonal}} & 27\\
\texttt{\detokenize{signal_decomposition}} & 30\\
\texttt{\detokenize{projector_residual_reconstructs}} & 34\\
\texttt{\detokenize{no_reconstruction_from_initial_probabilities}} & 40\\
\texttt{\detokenize{recovered_probability_at}} & 53\\
\texttt{\detokenize{projector_not_injective}} & 60\\
\texttt{\detokenize{publication_bridge}} & 68\\
\bottomrule\end{longtable}\normalsize

\section{What completion means here}
The fixed twelve-dimensional bridge is formally covered in the preserved later
checkpoint. This guide and machine concordance expose its proof and trust
boundaries.

General-modulus M\"obius and divisor-projector theory, polyphase identities, and a
general measurement API remain outside this Lean package. The traditional paper
supplies separate proofs. The current pinned replay passed for the fixed formal package.


\section{The original finite bridge and its extension}
\hypertarget{q:formal-coverage-lineage-and-reproducibility}{%
\section{Formal coverage, lineage, and
reproducibility}\label{q:formal-coverage-lineage-and-reproducibility}}

\hypertarget{q:two-inherited-checkpoints-and-the-expanded-boundary}{%
\subsection{Two inherited checkpoints and the expanded
boundary}\label{q:two-inherited-checkpoints-and-the-expanded-boundary}}

The earlier checked package has four mathematical production modules and
fifty-nine named theorem declarations. Its first module preserves
twenty-two original statements. Its proof-boundary document correctly
says that explicit matrix transport between independently defined
representations and normalized state-action identities are not supplied
as additional theorems in that package. Those limitations must not be
erased when describing the earlier evidence.

\Needspace{16\baselineskip}
The later checked package has five production modules and ninety-six
named theorem declarations:

{\small\begin{longtable}[]{@{}
  >{\raggedright\arraybackslash}p{\dimexpr 0.3333\columnwidth - 1.3332\tabcolsep\relax}
  >{\raggedleft\arraybackslash}p{\dimexpr 0.3333\columnwidth - 1.3332\tabcolsep\relax}
  >{\raggedright\arraybackslash}p{\dimexpr 0.3333\columnwidth - 1.3332\tabcolsep\relax}@{}}
\toprule\noalign{}
\begin{minipage}[b]{\linewidth}\raggedright
Module
\end{minipage} & \begin{minipage}[b]{\linewidth}\raggedleft
Named theorems
\end{minipage} & \begin{minipage}[b]{\linewidth}\raggedright
Principal role
\end{minipage} \\
\midrule\noalign{}
\endhead
\bottomrule\noalign{}
\endlastfoot
\texttt{QuantumArithmeticBridge} & 22 & Literal rational projector and
cyclotomic tables \\
\texttt{FiniteFourier} & 17 & Indexed exact Fourier and normalized state
identities \\
\texttt{SemanticBridge} & 30 & Sound complex interpretation and
transported identities \\
\texttt{Conventions} & 17 & Exponential entries, scales, and complex
projector laws \\
\texttt{Publication} & 10 & Phase lookup, residual reconstruction,
nonreconstruction, and central bundle \\
Total & 96 & Finite modulus-twelve proof chain, including helpers \\
\end{longtable}}

The later source explicitly supplies \(FDF^\dagger=P\), both normalized
state-action identities, the conventional exponential formula for \(F\),
identical initial probabilities, orthogonality, exact selected
probability one versus zero, histogram nonreconstruction, and projector
noninjectivity. It therefore closes the representation-transport
obligations that the earlier package leaves outside its formal boundary.

The ninety-six count includes helper lemmas and bundled propositions. It
is not a count of ninety-six independent discoveries. Nor should the
earlier fifty-nine and later ninety-six counts be added as if their
shared declarations were disjoint.

\hypertarget{q:current-and-inherited-checks}{%
\subsection{Current and inherited
checks}\label{q:current-and-inherited-checks}}

This edition preserves the source, theorem index, verification records,
and compiled-artifact lineage of the checked checkpoints. The archived
records provide the earlier compiler and kernel-checking evidence.
Current static audits and independent Python replay provide current
bounded evidence about the downloaded files and the example
implementation. These are disclosed separately in the release records.

The pinned runtime was restored for this completion pass. Fresh source compilation, axiom inspection, and replay of all five modules passed. The separate historical failures are preserved in the source records.

The exact source uses kernel-reduction decision proofs for finite
identities and ordinary proof terms for the transport lemmas. The
declaration audit rejects admitted proofs and unapproved
project-specific axioms. Standard foundational dependencies are reported
in the formal guide. Compiled artifact checking reuses Lean's kernel; it
is not an independently implemented proof assistant.

The later checkpoint preserves three negative controls with distinct
meanings. The false-sign claim is a directly false finite statement. The
false-mode case is an operational rejection whose failure mechanism must
be read alongside its diagnostic and the positive checked transform
table. The admitted-proof case tests the verification policy for
incomplete proof evidence. These controls are not combined with the four
controls of the earlier checkpoint into an artificial total.

\hypertarget{q:independent-arithmetic-evidence}{%
\subsection{Independent arithmetic
evidence}\label{q:independent-arithmetic-evidence}}

The original reproducibility package records 3,003 exact Python
assertions, separately from floating-point checks. Its exact verifier
represents the twelfth root in a cyclotomic quotient, checks the signal
and transform tables, the rational projector identities, the test-signal
decomposition, divisor-shell completeness, and the phase-sensitive
example. Its arithmetic regression range is bounded to moduli two
through sixty-four.

The release also includes the current replay and any added finite checks
with their own records. Exact rational or cyclotomic arithmetic and
numerical residuals are reported separately. A small floating-point
residual verifies the implementation at its declared tolerance; it does
not replace the exact proof of the general formula. Likewise, a finite
modulus sweep does not prove the arbitrary-modulus theorem. QA01 through
QA09 provide the general proofs needed here.

\hypertarget{q:the-twenty-two-preserved-original-statements}{%
\subsection{The twenty-two preserved original
statements}\label{q:the-twenty-two-preserved-original-statements}}

The original bridge module's statements remain available under their
original names. The following list records their substantive content
rather than treating each source helper as a separate main theorem.

{\small\begin{longtable}[]{@{}
  >{\raggedright\arraybackslash}p{\dimexpr 0.5000\columnwidth - 1.0000\tabcolsep\relax}
  >{\raggedright\arraybackslash}p{\dimexpr 0.5000\columnwidth - 1.0000\tabcolsep\relax}@{}}
\toprule\noalign{}
\begin{minipage}[b]{\linewidth}\raggedright
Original declaration
\end{minipage} & \begin{minipage}[b]{\linewidth}\raggedright
Preserved statement
\end{minipage} \\
\midrule\noalign{}
\endhead
\bottomrule\noalign{}
\endlastfoot
\texttt{ramanujan12\_length} & The rational signal list has length
twelve \\
\texttt{projector12\_shape} & The literal matrix has twelve rows and
twelve columns \\
\texttt{projector12\_symmetric} & The rational matrix equals its
transpose \\
\texttt{projector12\_idempotent} & Exact rational multiplication gives
\(P^2=P\) \\
\texttt{projector12\_fixes\_ramanujan} & The matrix fixes the signal
\(c\) \\
\texttt{projector12\_kills\_noise} & The matrix sends the displayed
perturbation \(b\) to zero \\
\texttt{projector12\_recovers\_test\_signal} & Applying it to \(x\)
returns \(c\) \\
\texttt{ramanujan12\_energy} & \(\|c\|^2=48\) \\
\texttt{noise12\_energy} & \(\|b\|^2=102\) \\
\texttt{testSignal12\_energy} & \(\|x\|^2=150\) \\
\texttt{energy\_decomposition} & The exact energy balance is
\(150=48+102\) \\
\texttt{zeta12\_order\_certificate} & The coordinate root raised to the
twelfth power is one \\
\texttt{zeta12\_inverse\_certificate} & The coordinate root times its
conjugate is one \\
\texttt{ramanujanC12\_table} & Four primitive phase sums give the exact
twelve-entry signal table \\
\texttt{ramanujan\_dft\_table} & The negative-sign DFT gives twelve
exactly at \(1,5,7,11\) \\
\texttt{abs\_ramanujan\_dft\_table} & The unsigned DFT gives the exact
even-support table \\
\texttt{same\_computational\_basis\_probabilities} & Signed and unsigned
rational entries have the same squares \\
\texttt{abs\_ramanujan12\_energy} & The unsigned signal also has squared
energy forty-eight \\
\texttt{recovered\_mode\_probability} & The scalar normalization gives
\(12^2/(48\cdot12)=1/4\) \\
\texttt{recovered\_total\_probability} & Four recovered mode
probabilities sum to one \\
\texttt{phase\_stripping\_coprime\_probability} & Zero transformed
coefficients give zero selected mass \\
\texttt{arithmetic\_signal\_quantum\_bridge} & The original finite
projector, recovery, DFT, and equal-square claims are bundled \\
\end{longtable}}

The later source strengthens the specification around these statements.
It proves that the four labels are exactly the coprime residues, that
the coordinate root has no smaller positive order, that the literal
rational and indexed matrices agree, and that the probability
consequences belong to the actual normalized complex state vectors.

The list operations in the original module use explicit finite literal
lists. Their generic helper definitions are not an assertion of
arbitrary-dimensional linear algebra for unsized lists, where zip
operations could truncate. General vector identities in this chapter use
fixed finite index sets, and the later matrix source uses
\texttt{Fin\ 12} for exactly that reason.

